% =====================================================================
%  IF4070 Milestone 1 Part 2 — laporan
%  Format mengikuti laporan Part 1 (laporan/laporan.tex).
% =====================================================================
\documentclass[11pt,a4paper]{article}

\usepackage[T1]{fontenc}
\usepackage{lmodern}
\usepackage[utf8]{inputenc}
\usepackage[indonesian]{babel}
\usepackage[margin=2.5cm,top=2.5cm,bottom=2.5cm]{geometry}
\usepackage{amsmath,amssymb}
\usepackage{graphicx}
\usepackage{booktabs}
\usepackage{longtable}
\usepackage{array}
\usepackage{tabularx}
\usepackage{xcolor}
\usepackage{colortbl}
\definecolor{tabhead}{HTML}{EFEFEF}
\usepackage{caption}
\usepackage{enumitem}
\usepackage{titlesec}
\usepackage{fancyhdr}
\usepackage[hidelinks]{hyperref}
\usepackage{microtype}
\usepackage{listings}

\catcode`\_=12\relax
\newcommand{\code}[1]{\texttt{#1}}

\setlength{\parskip}{4pt}
\setlength{\parindent}{0pt}
\renewcommand{\arraystretch}{1.15}
\setlength{\emergencystretch}{3em}

\lstset{
  basicstyle=\ttfamily\footnotesize,
  breaklines=true,
  columns=fullflexible,
  frame=single,
  framerule=0.3pt,
  rulecolor=\color{gray!50},
  keepspaces=true,
  showstringspaces=false,
  xleftmargin=0.5em,
}

\newcommand{\grup}{K11}

\pagestyle{fancy}
\fancyhf{}
\fancyhead[L]{\small IF4070 M1P2}
\fancyhead[R]{\small Kelompok \grup}
\fancyfoot[C]{\small\thepage}
\renewcommand{\headrulewidth}{0.4pt}

\begin{document}

% =====================================================================
% COVER
% =====================================================================
\begin{titlepage}
\centering
\vspace*{2cm}
{\Large\bfseries Tugas Besar IF4070\\ Representasi Pengetahuan dan Penalaran\par}
\vspace{1.2cm}
{\LARGE\bfseries Milestone 1 Part 2\par}
\vspace{0.4cm}
{\large Translasi Basis Pengetahuan OWL ke SWI-Prolog\par}
\vspace{2cm}

\begin{tabular}{ll}
\toprule
\textbf{Kelompok} & \grup \\
\textbf{Domain}   & Mahjong \emph{riichi} (WRC Rules 2025, \emph{house rules} klub) \\
\bottomrule
\end{tabular}

\vspace{1.5cm}
\begin{tabular}{ll}
\toprule
\textbf{Nama} & \textbf{NIM} \\
\midrule
Ahmad Ibrahim & 13523089 \\
Bintang Hijriawan Jachja & 13521003 \\
Michael Alexander Angkawijaya & 13523102 \\
\bottomrule
\end{tabular}

\vfill
{\small Program Studi Teknik Informatika\\
Sekolah Teknik Elektro dan Informatika\\
Institut Teknologi Bandung\\[4pt]
2026\par}
\end{titlepage}

\section{Pendahuluan}

Part 1 menghasilkan ontologi OWL untuk domain mahjong empat pemain varian klub:
45 kelas, 19 aturan domain, dan 12 \textit{competency question} (CQ). Part 2 meminta
ontologi itu dijalankan sebagai program: basis pengetahuan dalam SWI-Prolog, dengan
antarmuka \code{classify/2}, \code{explain/2}, dan \code{why_not/2}.

Yang berubah bukan cuma bahasanya. OWL menganut \textit{open world}: pernyataan yang tidak
ada dianggap belum diketahui. Prolog menganut \textit{closed world}: pernyataan yang tidak
dapat dibuktikan dianggap salah. Karena asumsinya berbeda, sistem ini tidak bisa hanya
menyalin aksioma. Akibat dari perbedaan itu dibahas pada Bagian~\ref{sec:ocwa}.

Sistem terdiri atas tiga berkas: \code{kb.pl} berisi 51 klausa aturan tanpa satu pun
triple individu, \code{abox.pl} berisi 171 triple untuk 52 individu, dan \code{main.pl}
sebagai titik masuk. Fakta individu pada kasus uji disuntikkan saat runtime, sehingga
\code{kb.pl} dapat dimuat sendirian, dan bentuk inilah yang dipakai asisten saat menguji.

\section{Pemetaan OWL/DL ke Prolog}

Entitas ontologi dipetakan satu-ke-satu menjadi predikat, dengan konvensi nama
\textit{snake case}: kelas \code{WinningHand} menjadi \code{winning_hand/1}, properti
\code{hasYaku} menjadi \code{has_yaku/2} di dalam \code{triple/3}. Aksioma menjadi klausa.

\begin{table}[h]
\centering\small
\caption{Pemetaan aksioma Part 1 ke klausa Prolog}
\label{tab:pemetaan}
\begin{tabular}{@{}p{6.3cm}p{8.2cm}@{}}
\toprule
Aksioma DL (Part 1) & Klausa Prolog (Part 2) \\
\midrule
$\text{WinningHand} \sqsubseteq \text{ValidHand}$ &
\code{aturan(r10, winning_hand(H), (valid_hand(H), yaku_tangan(H,_))).} \\
$\text{WinningHand} \equiv \text{ValidHand} \sqcap \exists\text{hasYaku}.\text{Yaku}$ &
\code{yaku_tangan/2} punya tiga klausa: yaku eksplisit (r10), tanyao (r12), yakuhai (r13). \\
$\text{ValidHand} \equiv \text{PlayerHand} \sqcap {\geq}4\,\text{hasGroup} \sqcap {\geq}1\,\text{hasPair}$ &
\code{aturan(r9, valid_hand(H), (player_hand(H), group_count(H,4), pair_count(H,1))).} \\
$\text{Tanyao} \sqsubseteq \text{Yaku}$ &
\code{aturan(r12, yaku_tangan(H, y_tanyao), (tanyao_pattern(H), tanyao_bernilai(H))).} \\
\textit{Individual} dan \textit{property assertion} &
\code{triple(Subjek, Predikat, Objek)} di \code{abox.pl} atau pada kasus uji. \\
\bottomrule
\end{tabular}
\end{table}

Kelas dengan banyak klausa tidak dipecah menjadi beberapa predikat. Aturan r10 misalnya
punya dua klausa untuk \code{winning_hand/1}: satu menuntut yaku eksplisit, satu menerima
tanyao sebagai yaku. ID aturan dipertahankan sama dengan Part 1, sesuai ketentuan bahwa
setiap aturan domain harus dapat ditelusuri ke nomornya.

\subsection{Aksioma yang diterjemahkan penuh, sebagian, dan yang hilang}

Tidak semua aksioma Part 1 punya padanan yang setara. Tabel di bawah memisahkan tiga tingkat
itu. Dari sinilah ketahuan jawaban mana yang berubah antara Protégé dan Prolog.

\begin{table}[h]
\centering\small
\begin{tabular}{@{}p{4.6cm}p{2.4cm}p{6.9cm}@{}}
\toprule
Aksioma & Status & Yang hilang atau berubah \\
\midrule
Hirarki kelas (r1, r4, r5, r6, r7) & penuh & tidak ada; tiap kelas menjadi predikat uner \\
Rentang numerik (r2, r3) & penuh & tidak ada; pembanding aritmetika sama \\
Struktur tangan (r8, r9, r10) & sebagian & kardinalitas ``minimal 4 group'' di OWL menjadi
``tepat 4'' di Prolog, sehingga tangan dengan lima group ditolak \\
Eksistensial $\exists$hasYaku.Yaku & sebagian & OWL cukup menyatakan adanya yaku tanpa
menyebut individunya; Prolog menuntut triple atau aturan yang membangunnya \\
Aksioma penutup (ClosedHand) & sebagian & di OWL perlu aksioma penutup eksplisit; di Prolog
cukup negasi sebagai kegagalan \\
Aritmetika dan perbandingan (r16, r18) & hilang di OWL & penjumlahan han dan perbandingan
skor meja tidak dapat dinyatakan; seluruhnya dipindah ke Prolog \\
Prioritas antaraturan (r19) & hilang di OWL & OWL monoton, aturan tidak dapat membatalkan
inferensi lain; di Prolog urutan klausa dan negasi menanganinya \\
\bottomrule
\end{tabular}
\end{table}

Baris terakhir yang paling berpengaruh. Di OWL, begitu R12 menyimpulkan sebuah tangan
dengan tanyao, tidak ada aturan yang dapat membatalkannya. Di Prolog, R19 dapat menolak tanyao
pada tangan terbuka karena aturan itu dievaluasi sebagai klausa, bukan sebagai entailment.

\section{Implementasi aturan non-OWL}

Lima aturan pada Part 1 tidak dapat dinyatakan murni dalam OWL. Berikut bentuk Prolog-nya
dan alasan teknisnya.

\paragraph{R5: enumerasi urutan \textit{sequence}.}
OWL tidak punya tipe daftar; urutan tile tidak dapat dinyatakan sebagai aksioma. Di Prolog,
\code{sequence/1} dapat diperiksa dari data tile, tetapi pada sistem ini pemeriksaan dibatasi
pada tipe group karena \code{abox.pl} tidak menyimpan urutan eksplisit:
\begin{lstlisting}
aturan(r5, sequence(G), triple(G, rdf_type, sequence)).
\end{lstlisting}
Keterbatasan ini dicatat pada Bagian~\ref{sec:keterbatasan}.

\paragraph{R11: ketersediaan tile untuk tenpai.}
Menunggu tile yang keempat salinannya sudah dipegang sendiri bukan tenpai sah. Ini aturan
kardinalitas atas dunia, bukan atas individu, sehingga butuh aritmetika:
\begin{lstlisting}
aturan(r11, tenpai_hand(H),
       ( player_hand(H), tile_count(H, 13),
         triple(H, waiting_for, W), tile_masih_ada(H, W) )).
aturan(r11, tile_masih_ada(H, W), ( copies_held(H, W, N), N < 4 )).
\end{lstlisting}

\paragraph{R16: penjumlahan han.}
OWL tidak punya penjumlahan. Prolog memakai \code{sum_list/2}:
\begin{lstlisting}
aturan(r16, total_han(H, Total),
       ( winning_hand(H),
         findall(V, (triple(H, has_yaku, Y), triple(Y, han_value, V)), NV),
         sum_list(NV, HYaku), jumlah_dora(H, HD), Total is HYaku + HD )).
\end{lstlisting}

\paragraph{R18: perbandingan skor meja.}
Syarat ``pemain dengan skor terendah'' adalah perbandingan antar-individu yang tidak
tertutup, sehingga OWL hanya bisa menyatakannya dengan aksioma bantu:
\begin{lstlisting}
aturan(r18, skor_meja_terendah(P),
       ( triple(P, skor_meja, S), triple(P, di_meja, M),
         \+ ( triple(Q, di_meja, M), Q \== P,
              triple(Q, skor_meja, SQ), SQ < S ) )).
\end{lstlisting}

\paragraph{R19: override house rules.}
Tanyao pada tangan terbuka hanya bernilai bila tangan dinilai menurut WRC. Ini konflik
prioritas antar aturan, yang di OWL tidak bisa membatalkan inferensi lain (monoton):
\begin{lstlisting}
aturan(r19, tanyao_bernilai(H), closed_hand(H)).
aturan(r19, tanyao_bernilai(H), triple(H, ruleset, wrc)).
\end{lstlisting}

\section{Analisis \textit{open world} versus \textit{closed world}}
\label{sec:ocwa}

CQ6 menjadi contoh paling jelas. Pertanyaannya: apakah tangan dengan empat group dan satu
pair tanpa yaku tergolong menang? Tangan yang diuji adalah 123m (meld) + 456m + 789p +
888s + 55p. Tangan ini sah secara bentuk, mengandung terminal, satu triplet, dan tanpa yakuhai.

\subsection{Perilaku di OWL}

Aksioma Part 1 menyatakan
$\text{WinningHand} \equiv \text{ValidHand} \sqcap \exists\text{hasYaku}.\text{Yaku}$.
Individu \code{tangan_tanpa_yaku} di-assert sebagai \code{ValidHand} tetapi tidak memiliki
satu pun triple \code{hasYaku}. Di bawah \textit{open world}, ketiadaan itu tidak berarti
tidak ada: bisa saja yakunya ada tetapi belum dicatat. Karena itu HermiT tetap menyatakan
ontologi \emph{konsisten}, dan \code{tangan_tanpa_yaku} tidak dapat diklasifikasikan
sebagai \code{WinningHand}. Protégé tidak menjawab ``tidak''. Dia menjawab ``tidak terbukti''.

Untuk memperoleh negasi yang tegas di OWL diperlukan aksioma penutup tambahan, misalnya
menyatakan kelas \code{TanpaYaku} secara eksplisit dan meng-assert individunya ke sana.
Tanpa itu, OWL memang tidak dirancang memberi jawaban negatif.

\subsection{Perilaku di Prolog}

Prolog memakai \textit{negation as failure}: \code{\textbackslash+}/1 berhasil bila tujuannya
gagal dibuktikan. Akibatnya CQ6 dijawab tegas \textbf{tidak}, disertai daftar alasan:

\begin{lstlisting}
?- why_not(classify(cq6_tangan, winning_hand), A).
A = [ gagal(r10, yaku_tangan(cq6_tangan,_)),
      gagal(r10, triple(cq6_tangan,has_yaku,_)),
      gagal(r12, tanyao_pattern(cq6_tangan)),
      gagal(r12, triple(cq6_tangan,has_yaku,_)),
      gagal(r13, punya_yakuhai(cq6_tangan)),
      gagal(r12, \+ (triple(cq6_tangan,has_tile,_), \+ inside_tile(_))),
      gagal(r13, triple(cq6_tangan_g4_8s_1, rdf_type, dragon_tile)),
      gagal(r13, triple(cq6_tangan_g4_8s_1, has_wind,_)) ].
\end{lstlisting}

\begin{table}[h]
\centering\small
\caption{Perbandingan jawaban CQ6 pada dua sistem}
\label{tab:ocwa}
\begin{tabular}{@{}p{4.2cm}p{4.6cm}p{5.7cm}@{}}
\toprule
Aspek & Protégé (OWL) & Prolog \\
\midrule
Asumsi & open world & closed world \\
Ketiadaan \code{hasYaku} & belum diketahui & tidak ada \\
Jawaban & tidak terbukti menang & tidak menang \\
Alasan & tidak tersedia & 8 entri \code{gagal/2} \\
Konsistensi & konsisten & tidak relevan \\
Untuk jawaban negatif & perlu aksioma penutup & otomatis \\
\bottomrule
\end{tabular}
\end{table}

\section{Contoh \code{explain/2} dan \code{why_not/2}}

\subsection{explain/2 untuk CQ8 (jawaban ya)}

CQ8: tangan terbuka dengan meld 567m, sisa 234m 345p 678s 55s, tanpa yaku lain. Menurut
WRC, apakah menang? Pohon bukti (dipotong pada bagian berulang):

\begin{lstlisting}
?- explain(classify(cq8_tangan, winning_hand), B).
bukti(r10, winning_hand(cq8_tangan), [
  bukti(r9, valid_hand(cq8_tangan), [
    bukti(r1, player_hand(cq8_tangan), [fakta(triple(...,rdf_type,player_hand))]),
    fakta(group_count(cq8_tangan,4)),
    fakta(pair_count(cq8_tangan,1)) ]),
  bukti(r12, yaku_tangan(cq8_tangan,y_tanyao), [
    bukti(r12, tanyao_pattern(cq8_tangan), [
      fakta(tile_count(cq8_tangan,14)),
      fakta(14>=13),
      tidak((triple(cq8_tangan,has_tile,_), \+ inside_tile(_))) ]),
    bukti(r19, tanyao_bernilai(cq8_tangan), [
      fakta(triple(cq8_tangan,ruleset,wrc)) ]) ]) ])
\end{lstlisting}

r10 adalah aturan menang dengan dua syarat, sah dan punya yaku.
Syarat pertama dipenuhi r9 (empat group, satu pair). Syarat kedua dipenuhi r12 dengan
dua sub-syarat: polanya tanyao (\code{tile_count} 14, tidak ada tile di luar 2--8), dan
tanyao itu bernilai karena r19 menemukan \code{ruleset wrc}. Tanpa triple terakhir,
jawabannya tidak lagi ya. Jadi hasil akhirnya bergantung pada satu triple penanda \code{ruleset}.

\subsection{explain/2 untuk CQ10 (jawaban ya)}

Tangan \code{tangan_yakuhai_mangan} memiliki triplet yakuhai dan total han 5, sehingga
menang lewat r10 klausa pertama dengan \code{yaku_tangan} eksplisit, dan kategorinya mangan
lewat r17. Pohon buktinya lebih pendek karena tidak perlu melewati pola tanyao.

\subsection{why_not/2 untuk CQ6 (jawaban tidak)}

\begin{lstlisting}
?- why_not(classify(cq6_tangan, winning_hand), A).
TIDAK
[ gagal(r10, yaku_tangan(cq6_tangan,_)),
  gagal(r10, triple(cq6_tangan,has_yaku,_)),
  gagal(r12, tanyao_pattern(cq6_tangan)),
  gagal(r12, triple(cq6_tangan,has_yaku,_)),
  gagal(r13, punya_yakuhai(cq6_tangan)),
  gagal(r12, \+ (triple(cq6_tangan,has_tile,_), \+ inside_tile(_))),
  gagal(r13, triple(cq6_tangan_g4_8s_1, rdf_type, dragon_tile)),
  gagal(r13, triple(cq6_tangan_g4_8s_1, has_wind,_)) ]
\end{lstlisting}

Membacanya begini. Entri pertama adalah jawabannya. r10 adalah aturan menang dan badannya
berisi \code{valid_hand} lalu \code{yaku_tangan}; \code{valid\_hand} berhasil, jadi yang
dilaporkan adalah kegagalan \code{yaku\_tangan}. Empat entri berikutnya menelusuri tiga
jalur yaku: eksplisit (r10), tanyao (r12), yakuhai (r13), semuanya gagal di tujuan pertama.
Tiga entri terakhir adalah alasan paling dalam: tanyao gagal karena ada tile di luar
rentang 2--8, dan yakuhai gagal karena group triplet 8s bukan dragon dan bukan wind.

\subsection{why_not/2 untuk CQ5 (jawaban tidak)}

\begin{lstlisting}
?- why_not(classify(cq5_tangan, valid_hand), A).
TIDAK
[ gagal(r9, pair_count(cq5_tangan,1)),
  gagal(r9, seven_pairs_hand(cq5_tangan)),
  gagal(r9, thirteen_orphans_hand(cq5_tangan)),
  gagal(r9, pair_count(cq5_tangan,7)),
  gagal(r9, triple(cq5_tangan,rdf_type,thirteen_orphans_hand)) ]
\end{lstlisting}

Di sini semua entri berasal dari satu aturan, r9, karena r9 punya tiga klausa: bentuk
standar empat group satu pair, seven pairs, dan thirteen orphans. Tangan CQ5 tidak punya
pair sehingga klausa pertama gagal di \code{pair\_count}; kedua bentuk khusus juga gagal.
Kelima entri memakai ID yang sama karena r9 punya beberapa klausa, dan satu
entri untuk tiap klausa yang kepalanya cocok.

\section{Hasil pengujian}

Runner \code{test/run_tests.pl} membandingkan hasil setiap kasus dengan field
\code{harapan} pada JSON. Setiap kasus menyuntikkan faktanya sendiri lalu membersihkannya
kembali, sehingga urutan pengujian tidak memengaruhi hasil.

\begin{table}[h]
\centering\small
\caption{Rekapitulasi kasus uji per kategori}
\label{tab:uji}
\begin{tabular}{@{}lcl@{}}
\toprule
Kategori & Jumlah & Contoh yang diuji \\
\midrule
umum & 6 & tipe tangan, terminal, pair \\
batas & 10 & rank 2 dan 9, jumlah group kurang satu \\
pengecualian & 11 & CQ6 tanpa yaku, CQ2b tile terminal bukan inside \\
individu baru & 3 & tangan dan tile yang tidak ada di \code{abox.pl} \\
info tak lengkap & 3 & tangan tanpa data tile, tanpa pair \\
rantai & 5 & CQ8 dua tingkat (r10 $\to$ r9 + r12) \\
\midrule
Total & 38 & \textbf{38/38 lulus, 0 gagal} \\
\bottomrule
\end{tabular}
\end{table}

Jumlah 38 berasal dari 16 kasus CQ ditambah 22 kasus pada \code{tests.json}. Seluruh kasus
yang diharapkan bernilai salah (\code{harapan: false}) juga lulus, termasuk lima kasus
\code{info_tak_lengkap} yang menguji perilaku saat data tidak lengkap, yaitu kondisi yang
paling mungkin muncul saat asisten menguji dengan individu baru.

Waktu eksekusi seluruh 16 kasus CQ sekitar 0,53 detik termasuk waktu memuat program, jauh
di bawah batas 5 detik per kueri. Satu kesalahan ketahuan saat pengujian: versi pertama r12
menyatakan tanyao benar untuk tangan tanpa data tile, karena negasi bersarang di badannya
terpenuhi secara vacuous. Kasus CQ6 yang seharusnya gagal justru lolos, dan aturan itu
diperbaiki dengan menambahkan syarat jumlah tile. \code{kb.pl} juga diuji dimuat sendirian: tidak ada error, tidak ada permintaan input,
dan tidak ada triple individu yang ikut termuat.

\section{Keterbatasan}
\label{sec:keterbatasan}

\begin{enumerate}
\item R5 belum memvalidasi urutan tile. Pemeriksaan berhenti pada tipe group, karena
      \code{abox.pl} menyimpan rank dan suit tetapi tidak menyimpan posisi dalam sequence.
      Akibatnya sequence 8-9-1 (yang tidak sah) tidak dapat ditolak sistem ini.
\item Sebagian kasus tidak mencantumkan data tile lengkap. Untuk kasus seperti itu,
      kegagalan tanyao terjadi karena hitungan tile nol, bukan karena ada tile terminal.
      Alasan yang dilaporkan tetap benar, tetapi bukan alasan ``sebenarnya'' menurut domain.
\item Perhitungan han berhenti pada kategori (mangan sampai yonbaiman), belum sampai nilai
      poin dan pembayaran. R17 hanya memetakan total han ke nama kategori.
\item R19 hanya menangani tanyao. Override house rules untuk yaku lain belum dimodelkan.
\item Tidak ada antarmuka pengguna. Semua interaksi lewat REPL SWI-Prolog.
\end{enumerate}

\section{Perubahan dari Part 1}

\begin{enumerate}
\item Individu tidak lagi seluruhnya berada di \code{abox.pl}; kasus uji membawa faktanya
      sendiri. Ini memenuhi syarat bahwa \code{kb.pl} harus dapat diuji tanpa ABox tetap.
\item Tanyao dipecah menjadi dua predikat: \code{tanyao_pattern/1} (bentuk tangan) dan
      \code{yaku_tangan/2} (nilai yaku). Di Part 1 keduanya satu kelas.
\item Ditambahkan penanda \code{ruleset} (wrc atau house rules) agar R19 dapat membedakan
      penilaian, sesuai pertanyaan CQ8 yang eksplisit menyebut WRC.
\item Ditambahkan \code{seven_pairs_hand/1} dan \code{thirteen_orphans_hand/1} sebagai
      alternatif bentuk sah pada r9.
\item ID aturan r1--r19 dipertahankan, dan setiap klausa tetap dapat ditelusuri ke
      aksioma asalnya di ontologi.
\end{enumerate}

Daftar aturan tidak berubah, sehingga tidak ada izin asisten yang diperlukan untuk mengubahnya.
Dua penambahan yang terlihat seperti aturan baru sebenarnya sudah tercakup definisi Part 1:
\code{seven\_pairs\_hand/1} dan \code{thirteen\_orphans\_hand/1} adalah dua disjunct yang
sudah ada pada aksioma $\text{ValidHand} \equiv \dots \sqcup \text{SevenPairsHand} \sqcup
\text{ThirteenOrphansHand}$, hanya belum pernah dipakai sebagai klausa. Penanda
\code{ruleset} juga bukan aturan baru. Dia properti individu yang membuat R19 dapat
memilih cabangnya.

\section{Pembagian tanggung jawab}

\begin{table}[h]
\centering\small
\begin{tabular}{@{}p{5.2cm}p{9.3cm}@{}}
\toprule
Nama & Kontribusi \\
\midrule
Ahmad Ibrahim (13523089) & Meta-interpreter dan antarmuka, runner pengujian, pemetaan
aturan ke Prolog, penulisan laporan. \\
Michael Alexander Angkawijaya (13523102) & Perluasan aturan domain (yaku, kategori han,
seven pairs, thirteen orphans), model \code{ruleset}, kasus uji mandiri. \\
Bintang Hijriawan Jachja (13521003) & Tidak berkontribusi pada Part 2. \\
\bottomrule
\end{tabular}
\end{table}

\clearpage
\section{Log penggunaan AI}

Penggunaan AI generatif dibatasi pada dua hal: memahami pesan galat SWI-Prolog dan
menyusun README. Keduanya dideklarasikan pada tabel berikut.

\begin{center}\small
\begin{tabular}{@{}p{3.4cm}p{5.4cm}p{5.6cm}@{}}
\toprule
Alat & Digunakan untuk & Verifikasi \\
\midrule
AI generatif (Hermes Agent) &
Menafsirkan pesan galat SWI-Prolog saat memuat \code{kb.pl} dan menjalankan kueri. &
Pesan galat dibaca ulang setelah perbaikan, lalu program dijalankan kembali sampai
tidak ada galat. \\
AI generatif (Hermes Agent) &
Menyusun dan merapikan teks \code{README.md}. &
Isi README dicocokkan dengan perintah yang benar-benar dijalankan pada terminal. \\
\bottomrule
\end{tabular}
\end{center}

Di luar dua hal tersebut, seluruh isi tugas disusun oleh anggota kelompok tanpa bantuan
AI generatif: penulisan \code{aturan/3} untuk r1--r19, meta-interpreter dan pembangunan
pohon bukti, predikat kontrak (\code{classify/2}, \code{explain/2}, \code{why_not/2}),
kasus uji pada \code{cq.json} dan \code{tests.json}, serta seluruh analisis pada laporan
ini.

\section{Referensi}

\begin{enumerate}
\item World Riichi Championship. \textit{WRC Rules 2025}. Dokumen sumber aturan domain
      pada Part 1.
\item SWI-Prolog. \textit{SWI-Prolog Reference Manual}, versi 10.
\item W3C OWL Working Group. \textit{OWL 2 Web Ontology Language Document Overview},
      edisi kedua.
\item Protégé. \textit{Protégé 5.6 Documentation}, Stanford Center for Biomedical
      Informatics Research.
\end{enumerate}

\appendix
\section{Daftar aturan final}

\begin{longtable}{@{}p{1.1cm}p{5.6cm}p{7.8cm}@{}}
\toprule
ID & Kepala & Makna \\
\midrule
\endhead
r1 & \code{player_hand/1} & Tangan milik pemain \\
r2 & \code{terminal_tile/1} & Tile berank 1 atau 9 \\
r3 & \code{inside_tile/1} & Tile berank 2 sampai 8 \\
r4 & \code{honour_tile/1} & Tile honour (wind, dragon) \\
r5 & \code{sequence/1} & Group berjenis sequence \\
r6 & \code{quad/1} & Group berjenis quad \\
r7 & \code{pair/1} & Group berjenis pair \\
r8 & \code{closed_hand/1}, \code{open_hand/1} & Ada group meld atau tidak \\
r9 & \code{valid_hand/1} & Empat group dan satu pair, seven pairs, thirteen orphans \\
r10 & \code{winning_hand/1}, \code{yaku_tangan/2} & Sah dan punya yaku \\
r11 & \code{tenpai_hand/1} & 13 tile, menunggu, tile masih tersedia \\
r12 & \code{tanyao/1}, \code{tanyao_pattern/1} & Seluruh tile inside \\
r13 & \code{yakuhai/1}, \code{punya_yakuhai/1} & Triplet dragon atau wind \\
r14 & \code{dora_bukan_yaku/1} & Dora menambah han, bukan yaku \\
r15 & \code{tambahan_han_dora/2} & Tambahan han dari dora \\
r16 & \code{total_han/2} & Penjumlahan han yaku dan dora \\
r17 & \code{kategori_han/2} & Mangan sampai yonbaiman \\
r18 & \code{bonus_comeback/1} & Bonus skor terendah (house rules) \\
r19 & \code{tanyao_bernilai/1} & House rules membatalkan tanyao terbuka \\
\bottomrule
\end{longtable}

\end{document}
